\section{Ограничения и интерпретация}
Гарантии предполагают точные градиенты и полную линейную минимизацию.
Неточные и стохастические оракулы требуют дополнительных слагаемых
ошибки и здесь не рассматриваются. Для логарифмического шага нужны
допустимые $L_0,L_1,D$, а для сертифицированной точности по функции
ещё и допустимый $\gamma$. Утверждение~\ref{prop:poisson} даёт эти
величины для конкретного класса, но не решает задачу их проверки
для произвольных моделей.

На компактном множестве наши предположения уже влекут обычную
гладкость. Существенна полезность параметров через зазор, а не
существование глобальной константы $L$. Утверждение~\ref{prop:rigidity}
ограничивает режим $L_0=0$: он не описывает непостоянную функцию с
нулевым градиентом в минимуме. Условие ненулевого градиента в
теореме~\ref{thm:strong_set_linear} содержательно.

Опыты выполнены на небольших синтетических задачах регрессии.
Общие и norm-based константы являются допустимыми аналитическими
границами, а не оптимизированными значениями. Время зависит от
оборудования и реализации. Результаты не доказывают всеобщего
превосходства, оптимальности зависимости от $\gamma$ или совпадения
верхних и нижних оценок сложности. Утверждение~\ref{prop:armijo}
прямо показывает, что известное адаптивное правило сохраняет тот же
порядок числа принятых обновлений. Вклад состоит в анализе кривизны
через зазор с явными параметрами и шаге, минимизирующем верхнюю модель
без отклонённых вычислений функции.

\paragraph{Вычислительные инструменты.}
ИИ использован при выводах, проверках, подготовке кода и текста.
За математические утверждения и результаты отвечают авторы.
Символьные тождества, проверки в арифметике с плавающей точкой и
математические доказательства разделены; формальная сертификация
системой проверки доказательств не заявляется.
